\chapter{Pengintegralan}\label{ch:g12:integ}

Pengintegralan menjawab dua pertanyaan sekaligus: berapa luas di bawah sebuah
kurva, dan bagaimana orang dapat memperoleh kembali sebuah \omterm{def:g11:func:function}{fungsi} dari laju
perubahannya? Teorema dasar kalkulus menyatakan bahwa keduanya pertanyaan yang
sama --- penemuan terdalam dan terpakai dari matematika abad ketujuh belas.

\section{Integral fungsi kontinu}

\begin{definition}[Integral fungsi positif]\label{def:g12:integ:area}
Misalkan $f$ \omterm{def:g12:limcont:continuity}{kontinu} dan positif pada $\intcc{a}{b}$. Yang disebut
\emph{integral}\index{integral}
\[
\int_a^b f(x)\,\dd x
\]
adalah luas, dalam satuan luas, daerah yang dibatasi oleh kurva $f$,
sumbu-$x$, dan garis tegak $x = a$ dan $x = b$.
\end{definition}

\begin{omfigure}
\begin{tikzpicture}
\begin{axis}[
  width=8.4cm, height=5cm,
  axis lines=middle, xmin=-0.4, xmax=3.4, ymin=-0.3, ymax=2.3,
  xtick={0.5,2.7}, xticklabels={$a$,$b$}, ytick=\empty,
  samples=100, domain=-0.2:3.3]
  \addplot[name path=curve, omDef, thick, domain=0.2:3.2]
    {1 + 0.5*sin(deg(2*x)) + 0.2*x};
  \path[name path=axis] (0.5,0) -- (2.7,0);
  \addplot[omDef!20] fill between[of=curve and axis, soft clip={domain=0.5:2.7}];
  \node at (axis cs:1.6,0.55) {$\displaystyle\int_a^b f$};
\end{axis}
\end{tikzpicture}
\end{omfigure}

Untuk \omterm{def:g11:func:function}{fungsi} yang tandanya sembarang, luas di bawah sumbu-$x$ dihitung
negatif; dan orang menetapkan
$\int_b^a f(x)\,\dd x = -\int_a^b f(x)\,\dd x$.

\begin{omfigure}
\begin{tikzpicture}
\begin{axis}[omaxis,
  width=9.6cm, height=4.8cm,
  xmin=-0.5, xmax=7, ymin=-1.4, ymax=1.4,
  xtick={3.1416,6.2832}, xticklabels={$\pi$,$2\pi$}, ytick={-1,1}]
  \addplot[name path=sine, omDef, thick, domain=0:6.2832, samples=120]
    {sin(deg(x))};
  \path[name path=xaxis] (0,0) -- (6.2832,0);
  \addplot[omDef!20] fill between[of=sine and xaxis,
    soft clip={domain=0:3.1416}];
  \addplot[omThm!20] fill between[of=sine and xaxis,
    soft clip={domain=3.1416:6.2832}];
  \node at (axis cs:1.57,0.4) {$+$};
  \node at (axis cs:4.71,-0.4) {$-$};
\end{axis}
\end{tikzpicture}

{\small Luas bertanda: $\int_0^{2\pi} \sin x\,\dd x = 0$, dengan daerah
di bawah sumbunya (merah) meniadakan daerah di atasnya (biru).}
\end{omfigure}

\begin{proposition}[Sifat integral]\label{prop:g12:integ:properties}
Misalkan $f, g$ \omterm{def:g12:limcont:continuity}{kontinu} pada suatu selang yang memuat $a, b, c$, dan
$\lambda \in \R$.
\begin{enumerate}
  \item \emph{Kelinearan:}
        $\displaystyle\int_a^b (f + \lambda g) = \int_a^b f + \lambda \int_a^b g$.
  \item \emph{Hubungan Chasles:}
        $\displaystyle\int_a^c f = \int_a^b f + \int_b^c f$.
  \item \emph{Kepositifan:} jika $a \leq b$ dan $f \geq 0$ pada $\intcc{a}{b}$,
        maka $\displaystyle\int_a^b f \geq 0$; jika $f \leq g$, maka
        $\displaystyle\int_a^b f \leq \int_a^b g$.
\end{enumerate}
\end{proposition}

\begin{proof}
Hubungan Chasles dan kepositifannya langsung terbaca dari tafsiran luasnya
(luas bertambah ketika daerahnya dijajarkan; daerah yang tingginya positif
berluas positif). Pembandingannya menyusul dengan menerapkan kepositifan pada
$g - f$. Kelinearannya jelas secara naluriah untuk jumlah \omterm{def:g11:func:function}{fungsi} positif
(tumpuklah luasnya) dan dibuktikan secara ketat di universitas;
\emph{di sini ia diterima saja}.
\end{proof}

\section{Teorema dasar kalkulus}

\begin{definition}[Primitif]\label{def:g12:integ:primitive}
Suatu \emph{primitif}\index{primitif} (atau antiturunan) dari $f$ pada selang
$I$ adalah \omterm{def:g11:func:function}{fungsi} $F$ yang \omterm{def:g12:deriv:derivative}{dapat diturunkan} pada $I$ dengan $F' = f$.
\end{definition}

\begin{proposition}\label{prop:g12:integ:primunique}
Jika $F$ suatu \omterm{def:g12:integ:primitive}{primitif} $f$ pada selang $I$, maka \omterm{def:g12:integ:primitive}{primitif} $f$ pada
$I$ tepat berupa \omterm{def:g11:func:function}{fungsi} $F + c$, $c \in \R$. Diberikan $x_0 \in I$ dan
$y_0 \in \R$, ada tepat satu \omterm{def:g12:integ:primitive}{primitif} dengan $F(x_0) = y_0$.
\end{proposition}

\begin{proof}
Jika $G' = F' = f$, maka $(G - F)' = 0$ pada selang $I$, sehingga $G - F$
bernilai tetap.\footnote{Bahwa \omterm{def:g11:func:function}{fungsi} yang \omterm{def:g12:deriv:derivative}{turunannya} nol pada suatu selang
bernilai tetap menyusul dari \cref{thm:g12:deriv:variations}: \omterm{def:g11:func:function}{fungsi} itu
sekaligus naik dan turun.} Sebaliknya setiap $F + c$ adalah \omterm{def:g12:integ:primitive}{primitif}. Syarat
$F(x_0) = y_0$ memakukan tetapannya.
\end{proof}

\begin{theorem}[Teorema dasar kalkulus]\label{thm:g12:integ:ftc}
\index{teorema dasar kalkulus}
Misalkan $f$ \omterm{def:g12:limcont:continuity}{kontinu} pada selang $I$ dan $a \in I$. \omterm{def:g11:func:function}{Fungsi}
\[
F \colon x \longmapsto \int_a^x f(t)\,\dd t
\]
adalah \omterm{def:g12:integ:primitive}{primitif} $f$ pada $I$ yang bernilai nol di $a$. Akibatnya, untuk
sebarang \omterm{def:g12:integ:primitive}{primitif} $G$ dari $f$ dan $a, b \in I$:
\[
\int_a^b f(t)\,\dd t = \bigl[G(t)\bigr]_a^b = G(b) - G(a).
\]
\end{theorem}

\begin{omfigure}
\begin{tikzpicture}
\begin{axis}[omaxis,
  width=9.2cm, height=5.6cm,
  xmin=-0.4, xmax=4.2, ymin=-0.4, ymax=2.9,
  xtick={0.5,2.2,2.7}, xticklabels={$a$,$x$,\ $x+h$}, ytick=\empty]
  \addplot[name path=curve, omDef, thick, domain=0:3.9, samples=80]
    {1 + 0.4*x + 0.2*sin(deg(2*x))};
  \path[name path=xaxis] (0,0) -- (3.9,0);
  \addplot[omDef!20] fill between[of=curve and xaxis,
    soft clip={domain=0.5:2.2}];
  \addplot[omProp!40] fill between[of=curve and xaxis,
    soft clip={domain=2.2:2.7}];
  \draw[black, dashed, thin] (axis cs:2.2,1.69) -- (axis cs:2.7,1.69);
  \draw[black, dashed, thin] (axis cs:2.2,1.925) -- (axis cs:2.7,1.925)
    node[above right, font=\small] {$f(x+h)$};
  \node at (axis cs:1.35,0.6) {$F(x)$};
\end{axis}
\end{tikzpicture}

{\small Gagasan pembuktiannya: pertambahan $F(x+h) - F(x)$ adalah luas
jalur sempit itu (jingga), yang terapit di antara persegi panjang bertinggi
$f(x)$ dan $f(x+h)$ di atas alas sepanjang $h$.}
\end{omfigure}

\begin{proof}[Bukti ketika $f$ naik]
Tetapkan $x \in I$ dan $h > 0$ dengan $x + h \in I$. Menurut Chasles,
\[
F(x+h) - F(x) = \int_x^{x+h} f(t)\,\dd t .
\]
Karena $f$ naik, berlaku $f(x) \leq f(t) \leq f(x+h)$ untuk
$t \in \intcc{x}{x+h}$, dan menurut pembandingan \omterm{def:g12:integ:area}{integral} (\omterm{def:g12:integ:area}{integral} \omterm{def:g11:func:function}{fungsi}
tetap $c$ pada selang sepanjang $h$ adalah $ch$):
\[
h\,f(x) \leq F(x+h) - F(x) \leq h\,f(x+h),
\]
sehingga
\[
f(x) \leq \frac{F(x+h) - F(x)}{h} \leq f(x+h).
\]
Ketika $h \to 0^+$, $f(x+h) \to f(x)$ menurut \omterm{def:g12:limcont:continuity}{kekontinuan}, dan teorema apit
memberi bahwa hasil bagi selisihnya menuju $f(x)$; kasus $h < 0$ setangkup
dengan itu. Jadi $F' = f$, dan $F(a) = 0$. Kasus \omterm{def:g12:limcont:continuity}{kontinu} umum, yang tidak
\omterm{def:g12:seq:monotonic}{monoton}, dibuktikan di universitas.

Akhirnya, jika $G$ sebarang \omterm{def:g12:integ:primitive}{primitif}, maka $G = F + c$
(\cref{prop:g12:integ:primunique}), sehingga
$G(b) - G(a) = F(b) - F(a) = \int_a^b f$.
\end{proof}

\begin{example}
$\displaystyle\int_0^1 x^2\,\dd x = \left[\frac{x^3}{3}\right]_0^1 =
\frac13$: luas di bawah \omterm{def:g11:quad:parabola}{parabolanya} sepertiga persegi satuan, seperti
yang sudah diketahui Archimedes.
\end{example}

Tabel \omterm{def:g12:integ:primitive}{primitif} dibaca dari tabel \omterm{def:g12:deriv:derivative}{turunannya}
($u$ menyatakan \omterm{def:g11:func:function}{fungsi} yang \omterm{def:g12:deriv:derivative}{dapat diturunkan}, $c$ suatu tetapan sebarang):
\begin{center}
\renewcommand{\arraystretch}{1.7}
\begin{tabular}{c|c||c|c}
$f(x)$ & \omterm{def:g12:integ:primitive}{primitif} & $f$ & \omterm{def:g12:integ:primitive}{primitif} \\
\hline
$x^n \ (n \neq -1)$ & $\dfrac{x^{n+1}}{n+1} + c$ &
$u'u^n \ (n \neq -1)$ & $\dfrac{u^{n+1}}{n+1} + c$\\
$\dfrac1x \ (x > 0)$ & $\ln x + c$ &
$\dfrac{u'}{u} \ (u > 0)$ & $\ln u + c$\\
$\eu^x$ & $\eu^x + c$ &
$u'\eu^u$ & $\eu^u + c$\\
$\cos x$ & $\sin x + c$ &
$\dfrac{u'}{\sqrt u} \ (u>0)$ & $2\sqrt u + c$\\
$\sin x$ & $-\cos x + c$ & &
\end{tabular}
\end{center}

\begin{method}[Mengenali bentuk $u' \times (\text{sesuatu dalam } u)$]\label{met:g12:integ:uprime}
Untuk mengintegralkan hasil kali, carilah faktor yang merupakan \omterm{def:g12:deriv:derivative}{turunan} \omterm{def:g11:func:function}{fungsi}
dalam $u$, sampai pada suatu tetapan pengali. Misalnya pada
$\int_0^1 x\,\eu^{x^2}\dd x$, faktor $x$ adalah $\frac12 (x^2)'$:
\[
\int_0^1 x\,\eu^{x^2}\dd x = \frac12\left[\eu^{x^2}\right]_0^1
= \frac{\eu - 1}{2}.
\]
\end{method}

\begin{theorem}[Pengintegralan parsial]\label{thm:g12:integ:ibp}
\index{pengintegralan parsial}
Misalkan $u, v$ \omterm{def:g12:deriv:derivative}{dapat diturunkan} pada $\intcc{a}{b}$ dengan \omterm{def:g12:deriv:derivative}{turunan} yang \omterm{def:g12:limcont:continuity}{kontinu}.
Maka
\[
\int_a^b u'(t)\,v(t)\,\dd t
= \bigl[u(t)\,v(t)\bigr]_a^b - \int_a^b u(t)\,v'(t)\,\dd t .
\]
\end{theorem}

\begin{proof}
Aturan hasil kali memberi $(uv)' = u'v + uv'$; dengan mengintegralkan kedua
ruasnya pada $\intcc{a}{b}$ dan memakai teorema dasarnya pada ruas kiri
diperoleh $\bigl[uv\bigr]_a^b = \int_a^b u'v + \int_a^b uv'$.
\end{proof}

\begin{example}\label{ex:g12:integ:ibp}
$\displaystyle\int_0^1 t\,\eu^{t}\,\dd t$: ambil $u' = \eu^t$, $v = t$, jadi
$u = \eu^t$, $v' = 1$:
\[
\int_0^1 t\,\eu^t\,\dd t = \bigl[t\,\eu^t\bigr]_0^1 - \int_0^1 \eu^t\,\dd t
= \eu - (\eu - 1) = 1 .
\]
\end{example}

\section{Penerapan}

\begin{definition}[Nilai rata-rata]\label{def:g12:integ:mean}
\emph{Nilai rata-rata}\index{nilai rata-rata} dari \omterm{def:g11:func:function}{fungsi} \omterm{def:g12:limcont:continuity}{kontinu} $f$ pada
$\intcc{a}{b}$ ($a<b$) adalah
\[
\mu = \frac{1}{b-a}\int_a^b f(t)\,\dd t .
\]
\end{definition}

\begin{proposition}\label{prop:g12:integ:meanbounds}
Jika $m \leq f \leq M$ pada $\intcc{a}{b}$, maka $m \leq \mu \leq M$.
\end{proposition}

\begin{proof}
Integralkan ketaksamaan $m \leq f(t) \leq M$ pada $\intcc{a}{b}$ lalu
bagilah dengan $b - a > 0$.
\end{proof}

\begin{method}[Luas di antara dua kurva]\label{met:g12:integ:areabetween}
Jika $f \geq g$ pada $\intcc{a}{b}$, luas di antara kedua kurvanya adalah
$\int_a^b \bigl(f(x) - g(x)\bigr)\dd x$. Jika kurvanya berpotongan, penggallah
selangnya di titik potongnya lalu integralkan $\abs{f - g}$ sepotong demi sepotong.
\end{method}

\begin{omfigure}
\begin{tikzpicture}
\begin{axis}[omaxis,
  width=8.4cm, height=5.8cm,
  xmin=-2.1, xmax=2.9, ymin=-1.8, ymax=5.6,
  xtick={-1,2}, ytick=\empty]
  \addplot[name path=line, omThm, thick, domain=-1.7:2.5] {x + 1};
  \addplot[name path=parab, omDef, thick, domain=-1.9:2.45] {x^2 - 1};
  \addplot[omDef!20] fill between[of=line and parab,
    soft clip={domain=-1:2}];
  \fill (axis cs:-1,0) circle (1.5pt);
  \fill (axis cs:2,3) circle (1.5pt);
\end{axis}
\end{tikzpicture}

{\small Luas di antara garis $y = x + 1$ (merah) dan \omterm{def:g11:quad:parabola}{parabola}
$y = x^2 - 1$ (biru) adalah $\int_{-1}^{2}
\bigl((x+1) - (x^2-1)\bigr)\dd x = \frac92$.}
\end{omfigure}

\section{Latihan}

\begin{exercise}[$\star$]\label{exo:g12:integ:1}
Hitunglah
\[
\int_1^2 \left(3x^2 - \frac{1}{x^2}\right)\dd x, \qquad
\int_0^{\pi/2} \cos t \,\dd t, \qquad
\int_0^{1} \frac{\dd t}{2t+1} .
\]
\end{exercise}

\begin{exercise}[$\star$]\label{exo:g12:integ:2}
Carilah \omterm{def:g12:integ:primitive}{primitif} $F$ dari $f(x) = x\eu^{x^2}$ pada $\R$ dengan $F(0) = 1$.
\end{exercise}

\begin{exercise}[$\star$]\label{exo:g12:integ:3}
Hitunglah \omterm{def:g12:integ:mean}{nilai rata-rata} $f(t) = \sin t$ pada $\intcc{0}{\pi}$, lalu
tafsirkan hasilnya pada sebuah \omterm{def:g10:functions:graph}{grafik}.
\end{exercise}

\begin{exercise}[$\star\star$]\label{exo:g12:integ:4}
Dengan \omterm{thm:g12:integ:ibp}{pengintegralan parsial}, hitunglah
\[
\int_1^{\eu} \ln t\,\dd t
\qquad\text{dan}\qquad
\int_0^{\pi} t \sin t\,\dd t .
\]
\end{exercise}

\begin{exercise}[$\star\star$]\label{exo:g12:integ:5}
Hitunglah luas daerah di antara \omterm{def:g11:quad:parabola}{parabola} $y = x^2$ dan garis
$y = x + 2$.
\end{exercise}

\begin{exercise}[$\star\star$]\label{exo:g12:integ:6}
Misalkan $I = \displaystyle\int_0^1 \frac{\dd t}{1 + t}$.
\begin{enumerate}
  \item Hitunglah $I$.
  \item Untuk $n \in \N$, misalkan $I_n = \displaystyle\int_0^1 \frac{t^n}{1+t}\dd t$.
        Tunjukkan bahwa $I_n + I_{n+1} = \dfrac{1}{n+1}$, dan bahwa
        $0 \leq I_n \leq \dfrac{1}{n+1}$.
  \item Simpulkan bahwa
        $1 - \frac12 + \frac13 - \dots + \frac{(-1)^{n-1}}{n}
        \xrightarrow[n \to +\infty]{} \ln 2$.
\end{enumerate}
\end{exercise}

\begin{exercise}[$\star\star$]\label{exo:g12:integ:7}
Laju sebuah kereta (dalam m/s) selama $100$ sekon pertama sesudah berangkat
dimodelkan oleh $v(t) = 30\bigl(1 - \eu^{-t/50}\bigr)$. Hitunglah jarak yang
ditempuhnya selama $100$ sekon itu, dan laju rata-rata keretanya pada
selang tersebut.
\end{exercise}

\begin{exercise}[$\star\star\star$]\label{exo:g12:integ:8}
Untuk $n \geq 1$, misalkan $S_n = \dfrac1n \displaystyle\sum_{k=1}^{n} \frac{1}{1 + k/n}$.
\begin{enumerate}
  \item Tafsirkan $S_n$ sebagai luas persegi panjang yang menghampiri daerah
        di bawah kurva $t \mapsto \frac{1}{1+t}$ pada $\intcc{0}{1}$.
  \item Dengan memakai kemonotonan $t \mapsto \frac{1}{1+t}$, tunjukkan bahwa
        \[
        S_n \leq \int_0^1 \frac{\dd t}{1+t} \leq S_n + \frac1n\left(1 - \frac12\right),
        \]
        lalu simpulkan $\lim\limits_{n\to+\infty} S_n = \ln 2$.
\end{enumerate}
\end{exercise}

\begin{exercise}[$\star\star\star$]\label{exo:g12:integ:9}
\emph{(\omterm{def:g12:integ:area}{Integral} Wallis.)} Untuk $n \in \N$, misalkan
$W_n = \displaystyle\int_0^{\pi/2} \sin^n t\,\dd t$.
\begin{enumerate}
  \item Hitunglah $W_0$ dan $W_1$.
  \item Dengan menulis $\sin^{n+2}t = \sin t \cdot \sin^{n+1} t$ lalu
        mengintegralkan secara parsial, tunjukkan bahwa $W_{n+2} = \dfrac{n+1}{n+2}\,W_n$.
  \item Simpulkan $W_2$, $W_3$, $W_4$ lalu tunjukkan bahwa $(W_n)$ turun dan
        positif.
\end{enumerate}
\end{exercise}

\section{Soal: Archimedes melawan mesin}

\begin{problem}[{Soal akhir pekan --- luas di bawah parabola,
dihitung dengan tiga cara sepanjang dua puluh dua abad, dan
jati diri rahasia logaritma sebagai sebuah luas}]
\label{pb:g12:integ:1}
Sekitar tahun 240 SM, Archimedes menghitung luas eksak sebuah tembereng
\omterm{def:g11:quad:parabola}{parabola} --- tanpa \omterm{def:g10:coordgeom:system}{koordinat}, tanpa limit, tanpa aljabar.
Sembilan belas abad kemudian, persegi panjang Riemann mengulanginya dengan
apitan yang kasar; dan teorema dasar kalkulus
(\cref{thm:g12:integ:ftc}) kini melakukannya dalam satu baris. Soal ini
memainkan ketiga pertandingan itu, lalu memakai mesin yang sama untuk
menyingkap apa sebenarnya \omterm{def:g12:exp:ln}{logaritma} itu: luas dengan kesetangkupan penskalaan.

\medskip
\noindent\textbf{Bagian I --- Kelancaran.}
\begin{enumerate}
  \item Hitunglah $\displaystyle\int_0^1 (3x^2 - 2x + 1)\,\dd x$,
        \ $\displaystyle\int_1^{\eu} \frac{\dd x}{x}$, \ dan
        $\displaystyle\int_0^{\pi/2} \cos x\,\dd x$.
  \item Kenalilah bentuk $u'u$
        (\cref{met:g12:integ:uprime}):
        $\displaystyle\int_0^1 x\,\eu^{x^2}\dd x$ dan
        $\displaystyle\int_0^1 \frac{2x}{x^2 + 1}\,\dd x$.
  \item Secara parsial (\cref{thm:g12:integ:ibp}):
        $\displaystyle\int_0^1 x\,\eu^x \dd x$ dan
        $\displaystyle\int_1^{\eu} \ln x\,\dd x$.
  \item Hitunglah \omterm{def:g12:integ:mean}{nilai rata-rata}
        (\cref{def:g12:integ:mean}) dari $\sin$ pada
        $\intcc{0}{\pi}$ --- lalu perhatikan bahwa nilainya \emph{bukan}
        $\frac12$.
  \item Hitunglah luas di antara garis $y = x$ dan
        \omterm{def:g11:quad:parabola}{parabola} $y = x^2$ pada $\intcc{0}{1}$
        (\cref{met:g12:integ:areabetween}).
\end{enumerate}

\medskip
\noindent\textbf{Bagian II --- \omterm{def:g11:quad:parabola}{Parabolanya} dengan tiga cara.}
\begin{enumerate}[resume]
  \item Rancangan Riemann untuk luas di bawah $y = x^2$ pada
        $\intcc{0}{1}$: tulislah jumlah bawah $L_n$ dan jumlah
        atas $U_n$ atas $n$ persegi panjang yang sama lebarnya (seperti pada
        \cref{exo:g12:integ:8}).
  \item Buktikan dengan induksi (mesin \cref{pb:g12:seq:1}) rumus
        jumlah kuadratnya
        \[
        1^2 + 2^2 + \dots + n^2 = \frac{n(n+1)(2n+1)}{6} .
        \]
  \item Simpulkan bentuk tertutup $U_n$ dan $L_n$, hitunglah
        limit persekutuannya, lalu simpulkan: luasnya $\frac13$.
  \item Sekarang mesinnya: hitunglah
        $\int_0^1 x^2\,\dd x$ dengan teorema dasar itu, dalam
        satu baris saja. Bandingkanlah jerih payahnya.
  \item Archimedes menyatakannya dengan cara lain: \emph{sebuah tembereng
        \omterm{def:g11:quad:parabola}{parabola} adalah $\frac43$ segitiga yang termuat di dalamnya}. Untuk
        tembereng yang dipotong dari $y = x^2$ oleh tali busur penghubung
        $(-1, 1)$ ke $(1, 1)$: hitunglah luas temberengnya dengan
        \omterm{def:g12:integ:area}{integral}, luas segitiga di dalamnya (puncaknya di
        titik puncak $(0, 0)$), lalu periksalah perbandingan sang guru itu.
  \item Cara Archimedes sendiri: isilah temberengnya dengan segitiga
        besar $T$, lalu dua segitiga yang berjumlah
        $\frac T4$, lalu empat yang berjumlah $\frac{T}{16}$, dan
        seterusnya. Jumlahkan deret ukurnya lalu peroleh kembali
        $\frac43 T$ --- batang cokelat yang digigit tanpa henti di
        jilid sekolah menengah pertama, yang dimakan oleh ahli ukur Yunani.
  \item Satu kalimat untuk masing-masingnya: penghabisan (Archimedes), apitan
        (Riemann), pengantiturunan (Newton--Leibniz) ---
        apa yang diperlukan masing-masing, dan apa yang diberikannya?
\end{enumerate}

\medskip
\noindent\textbf{Bagian III --- \omterm{def:g12:exp:ln}{Logaritma} adalah sebuah luas.}
Untuk $x > 0$, ambil
$A(x) = \displaystyle\int_1^x \frac{\dd t}{t}$.
\begin{enumerate}[resume]
  \item Berikan $A(1)$ dan $A'(x)$
        (\cref{thm:g12:integ:ftc}), lalu simpulkan bahwa $A$
        tepat merupakan \omterm{def:g12:exp:ln}{logaritma asli} pada
        \cref{def:g12:exp:ln}.
  \item Keajaiban penskalaannya: tetapkan $a > 0$ lalu kajilah
        $g(x) = A(ax) - A(x)$. Hitunglah $g'$ (aturan rantai),
        simpulkan bahwa $g$ bernilai tetap, hitunglah tetapannya ---
        lalu simpulkan \omterm{def:g10:algebra:equation}{persamaan} fungsionalnya
        \[
        \ln(ab) = \ln a + \ln b ,
        \]
        yang terbukti dengan kalkulus murni: luas dari $1$ sampai $ab$
        terpecah menjadi salinan-salinan yang terskala.
  \item Simpulkan dari pertanyaan 14: $\ln(a^n) = n\ln a$ dan
        $\ln\frac1a = -\ln a$.
  \item \Cref{exo:g12:integ:8} membaca
        $\frac{1}{n+1} + \frac{1}{n+2} + \dots + \frac{1}{2n}$
        sebagai persegi panjang di bawah $\frac{1}{1 + t}$: hitunglah jumlah
        itu untuk $n = 10$ (tiga angka di belakang titik) lalu bandingkan dengan
        $\ln 2$. Jumlah bercita rasa harmonis yang mana, yang suku demi
        sukunya divergen, di sini justru \omterm{def:g12:seq:limit}{konvergen} ke sebuah luas?
\end{enumerate}

\medskip
\noindent\textbf{Bagian IV --- Penimbunan.}
\begin{enumerate}[resume]
  \item Sebuah mobil dipercepat dengan laju $v(t) = 3t^2$ m/s untuk
        $t \in \intcc{0}{10}$. Hitunglah jarak yang ditempuhnya,
        laju rata-ratanya, dan saat ketika
        laju sesaatnya sama dengan laju rata-ratanya.
  \item Sebuah batang sepanjang $4$ meter berkerapatan linear
        $\rho(x) = 2 + x$ kg/m. Hitunglah massa totalnya dan pusat
        massanya
        $\dfrac{\int_0^4 x\,\rho(x)\,\dd x}{\int_0^4
        \rho(x)\,\dd x}$ --- titik seimbang segitiga karton di
        jilid sekolah menengah pertama, yang akhirnya
        terhitung dengan bobot yang berubah-ubah.
  \item \omterm{def:g12:integ:area}{Integral} Wallis (\cref{exo:g12:integ:9}): hitunglah
        $W_0$, $W_1$, dan
        $W_2 = \int_0^{\pi/2} \sin^2 t\,\dd t$ dengan memakai
        pelinearan pada \cref{pb:g12:trigo:1}. (Tangga tak berhingganya
        mendaki, di jilid universitas, sampai ke rumus hasil kali
        untuk $\pi$ dan ke penormalan kurva
        loncengnya.)
  \item Penutup --- tiga wajah pengintegralan: sebuah \emph{luas}
        menurut definisinya, sebuah \emph{penimbunan} di dunia nyata
        (jarak, massa), sebuah \emph{antiturunan} menurut
        teorema dasarnya; dan sejarahnya dalam tiga nama.
        Tutuplah dengan mutiara bab ini: tombol sehari-hari yang mana
        pada kalkulator yang diam-diam merupakan luas di bawah
        $\frac1t$ --- dan identitas mana yang dibuktikan oleh luas itu?

\end{enumerate}
\end{problem}
